\documentclass[11pt]{scrartcl}
\usepackage[boxthm,secthm,noasy]{evan}
\usepackage{booktabs}
\renewcommand{\S}{\ensuremath{\mathsection}}
\renewcommand{\labelitemi}{\ensuremath{\bullet}}

\title{Week 1: Contradiction and Induction}
\author{Claude}
\date{2 November 2026}

\begin{document}
\maketitle
\begin{center}
  \textbf{Reading:} P\&B \S1.1, \S1.2 (pp.~1--7, up to strong induction).
  ACPS \S1.2, \S2.1--2.3.
\end{center}

\section{How to work on a problem}

\subsection{Problems and exercises}
An \vocab{exercise} is a question you know how to answer as soon as you read it.
A \vocab{problem} is a question you do not know how to start.
This course is about problems. On a problem, the first 20 minutes usually produce
no visible progress, and that is normal.

Every solution has two parts (ACPS \S2.2):
\begin{itemize}
  \ii the \vocab{investigation}, where you find out what is true and why; and
  \ii the \vocab{argument}, where you write a proof that convinces a reader.
\end{itemize}
The finished proof in a book shows only the argument. These notes show both,
because the investigation is the part you need to learn.
In each worked example below, the paragraph \textsf{Investigation} is the thinking
and the \textsf{Proof} is what you would hand in.

\subsection{Orientation}
Do this at the start of every problem.
\begin{enumerate}
  \ii Read the problem slowly. Note details: positive or nonnegative, integer or real,
  finite or infinite, strict or non-strict inequality.
  \ii Classify it: is it ``find'' or ``prove''? Have you seen a similar problem?
  \ii Write the \vocab{hypothesis} (what is given) and the \vocab{conclusion}
  (what you must show or find). Introduce notation for them.
  \ii Restate the conclusion in a different form if you can.
\end{enumerate}

\subsection{Four strategies for getting started}
When you are oriented but have no idea, try these (ACPS \S2.2).
\begin{description}
  \ii[Get your hands dirty.] Compute small cases. Make a table. Look for a pattern,
  then ask why the pattern happens.
  \ii[Penultimate step.] Ask: ``What fact would give the conclusion in one step?''
  Then try to prove that fact instead. For example, to show $A = B$, it is enough
  to show $A \le B$ and $B \le A$. To show a number is not a perfect square,
  it is enough to show it lies strictly between two consecutive squares.
  \ii[Wishful thinking.] Ask: ``What makes this problem hard?'' Then pretend that difficulty
  is not there and solve the problem without it. Then see how much of that solution survives
  when you put the difficulty back.
  \ii[Make it easier.] Solve a special case: smaller numbers, fewer variables,
  $n = 1, 2, 3$, an extra assumption. The special case often shows the idea of the general case.
\end{description}

\subsection{Strategy, tactic, tool}
ACPS uses three levels of method:
\begin{itemize}
  \ii a \vocab{strategy} is a general plan that applies to almost any problem
  (the four above, for example);
  \ii a \vocab{tactic} is a method that applies to many problems of one kind
  (induction, contradiction, pigeonhole, invariants);
  \ii a \vocab{tool} is a specific technique for a narrow task
  (``a polynomial with integer coefficients satisfies $a - b \mid P(a) - P(b)$'').
\end{itemize}
The \vocab{crux move} is the one step that makes the rest of the solution routine.
When you read a solution, find its crux move and write it down in one sentence.
That sentence is what you will remember the next time.

\begin{example}[ACPS 1.2.1, Hungary 1926]
  Prove that the product of four consecutive positive integers is never a perfect square.
\end{example}

\paragraph{Investigation.}
Orientation: hypothesis is ``$n$ is a positive integer''; conclusion is
``$n(n+1)(n+2)(n+3)$ is not a square.''

Penultimate step: how do you show a number is \emph{not} a square?
One way is to show it lies strictly between two consecutive squares $m^2$ and $(m+1)^2$,
because no square lies there. So we want to find $m$ depending on $n$.

Get your hands dirty:
\[
  1\cdot2\cdot3\cdot4 = 24,\quad 2\cdot3\cdot4\cdot5 = 120,\quad
  3\cdot 4\cdot 5\cdot 6 = 360,\quad 4\cdot5\cdot6\cdot7 = 840.
\]
These are $5^2 - 1$, $11^2 - 1$, $19^2 - 1$, $29^2 - 1$. So the product seems to be
\emph{one less than a square}. That is exactly the penultimate step we wanted:
a number of the form $s^2 - 1$ with $s \geq 2$ lies strictly between $(s-1)^2$ and $s^2$.

Why is it one less than a square? Pair the outer factors and the inner factors:
$n(n+3) = n^2 + 3n$ and $(n+1)(n+2) = n^2 + 3n + 2$. The two products differ by $2$.
With $m = n^2 + 3n$, the product is $m(m+2) = (m+1)^2 - 1$.
(Check: $n = 1$ gives $m = 4$ and $5^2 - 1 = 24$.)

\begin{proof}
Let $m = n^2 + 3n \geq 4$. Then
\[
  n(n+1)(n+2)(n+3) = (n^2+3n)(n^2+3n+2) = m(m+2) = (m+1)^2 - 1.
\]
Since $m^2 < m^2 + 2m = (m+1)^2 - 1 < (m+1)^2$, the product lies strictly between the
consecutive squares $m^2$ and $(m+1)^2$, so it is not a square.
\end{proof}

\begin{remark}
  The tool used here, \emph{squeeze between consecutive squares}, appears often.
  The strategies that found it were penultimate step (decide what kind of fact
  would finish the problem) and get your hands dirty (find the specific square).
\end{remark}

\section{Argument by contradiction}

\subsection{The method}
To prove a statement $S$, assume $S$ is false and deduce something impossible.
Then $S$ must be true.

\paragraph{When to use it.}
Use contradiction when the \emph{negation} of the conclusion gives you something
concrete to work with. Typical signs:
\begin{itemize}
  \ii the conclusion says something does \emph{not} exist or \emph{cannot} happen
  (``there is no polynomial \ldots'', ``is irrational'');
  then the negation gives you an object to study;
  \ii the conclusion says there are \emph{infinitely many} of something;
  then the negation gives you a finite list;
  \ii the conclusion says \emph{some} object has a property;
  then the negation says \emph{every} object fails it, which is a statement about all objects at once.
\end{itemize}
ACPS (\S2.3) puts it this way: ask ``What happens if I negate the conclusion?
Do I get something I can work with?''

\paragraph{Negating correctly.}
Most errors in contradiction proofs come from a wrong negation. The rules:
\begin{center}
\begin{tabular}{ll}
  \toprule
  \textbf{Statement} & \textbf{Negation} \\
  \midrule
  for every $x$, $Q(x)$ & there is an $x$ with not $Q(x)$ \\
  there is an $x$ with $Q(x)$ & for every $x$, not $Q(x)$ \\
  $A$ and $B$ & not $A$, or not $B$ \\
  if $A$ then $B$ & $A$ and not $B$ \\
  \bottomrule
\end{tabular}
\end{center}

\begin{theorem}[Euclid]
  There are infinitely many primes.
\end{theorem}

\paragraph{Investigation.}
The conclusion is ``infinitely many'', so the negation gives a finite list
$p_1, \dots, p_n$ of all primes. A finite list of numbers can be multiplied together.
The product $p_1 \cdots p_n$ is divisible by every prime, so $p_1 \cdots p_n + 1$
is divisible by none of them. But every integer greater than $1$ has a prime factor.

\begin{proof}
Assume, for contradiction, that $p_1, \dots, p_n$ are all the primes.
Let $N = p_1 p_2 \cdots p_n + 1 > 1$, and let $p$ be a prime factor of $N$.
Then $p = p_i$ for some $i$, so $p$ divides both $N$ and $p_1 \cdots p_n$,
hence $p$ divides their difference $1$. This is impossible.
\end{proof}

Before the next example, we need a tool that is used all through the course.

\begin{fact}\label{fact:poly}
  Let $P$ be a polynomial with integer coefficients and let $a, b$ be integers.
  Then $a - b$ divides $P(a) - P(b)$.
  In particular, if $a \equiv b \pmod m$, then $P(a) \equiv P(b) \pmod m$.
\end{fact}
\begin{proof}
If $P(x) = \sum c_j x^j$, then $P(a) - P(b) = \sum c_j (a^j - b^j)$, and each term is divisible by
$a - b$ because $a^j - b^j = (a-b)(a^{j-1} + a^{j-2}b + \dots + b^{j-1})$.
\end{proof}

\begin{example}[P\&B \S1.1, Euler]
  Prove that there is no polynomial $P$ with integer coefficients and degree at least $1$
  such that $P(0), P(1), P(2), \dots$ are all prime.
\end{example}

\paragraph{Investigation.}
The conclusion says ``there is no polynomial'', so assume there is one and study it.
Penultimate step: to get a contradiction, it is enough to find one $n \geq 0$ with $P(n)$ not prime.

Get your hands dirty with a famous near-example, $P(x) = x^2 + x + 41$, which is prime for
$x = 0, 1, \dots, 39$. Where does it fail? At $x = 40$: $40^2 + 40 + 41 = 41^2$.
And at $x = 41$: $41^2 + 41 + 41 = 41 \cdot 43$. The failure involves $41$, which is $P(0)$.
Try another: $P(x) = 2x + 3$. Then $P(0) = 3$ and $P(3) = 9 = 3 \cdot 3$. Again the value
$P(0)$ divides a later value.

So the guess is: $p = P(0)$ divides $P(p)$, $P(2p)$, $P(3p), \dots$.
This is \Cref{fact:poly} with $a = kp$, $b = 0$: $kp - 0$ divides $P(kp) - P(0)$,
so $p \mid P(kp)$.
If all values are prime, a prime divisible by $p$ must equal $p$. So $P$ takes the value $p$
infinitely often. A nonconstant polynomial cannot do that.

\begin{proof}
Assume, for contradiction, that such a $P$ exists. Let $p = P(0)$, a prime.
For each integer $k \geq 1$, \Cref{fact:poly} gives $kp \mid P(kp) - P(0)$, so $p \mid P(kp)$.
Since $P(kp)$ is a prime divisible by $p$, we get $P(kp) = p$.
Then the polynomial $P(x) - p$ has the infinitely many roots $p, 2p, 3p, \dots$,
so $P(x) - p$ is the zero polynomial and $P$ is constant. This contradicts $\deg P \geq 1$.
\end{proof}

\begin{remark}
  Two tools appeared here: \Cref{fact:poly}, and ``a nonzero polynomial has only finitely many roots''.
  The second one turns ``takes the same value infinitely often'' into ``is constant''.
\end{remark}

\begin{example}[P\&B \S1.1, Tomescu]
  Let $\SF = \{E_1, E_2, \dots, E_s\}$ be a family of subsets of a set $X$, each with exactly $r$
  elements. Suppose that any $r+1$ sets in $\SF$ (not necessarily distinct) have nonempty intersection.
  Prove that the intersection of all the sets in $\SF$ is nonempty.
\end{example}

\paragraph{Investigation.}
Make it easier: take $r = 1$. Each set is a single point, and any two of them intersect,
so any two are equal. Then all sets are the same point and the total intersection is that point.

For general $r$, negate the conclusion: the total intersection is empty.
That means: \emph{every} element $x$ is missing from \emph{some} set in $\SF$.
This is a statement about every element, which we can apply to chosen elements.

Penultimate step: to contradict the hypothesis, we need $r+1$ sets with empty intersection.
Which elements do those sets have to exclude? If one of the sets is $E_1$, the intersection
is already inside $E_1$, so we only need to exclude the $r$ elements of $E_1$.
Each of those $r$ elements is missing from some set. That gives $r$ more sets, and
$1 + r$ sets in total. The count matches the $r+1$ in the hypothesis exactly,
which is good evidence that this is the intended argument.

\begin{proof}
Assume, for contradiction, that $\bigcap_{j} E_j = \varnothing$.
Write $E_1 = \{x_1, \dots, x_r\}$. For each $i$, the element $x_i$ is not in every set of $\SF$,
so there is a set $E_{j_i} \in \SF$ with $x_i \notin E_{j_i}$.
The $r+1$ sets $E_1, E_{j_1}, \dots, E_{j_r}$ have intersection contained in $E_1$ and
containing none of $x_1, \dots, x_r$, so the intersection is empty.
This contradicts the hypothesis.
\end{proof}

\begin{remark}
  The numbers in a problem statement are information. Here ``$r+1$'' told us how many sets
  to build. If your argument never uses a number or a hypothesis from the statement,
  check whether the argument is wrong or the statement is stronger than needed.
\end{remark}

\subsection{Contradiction or contrapositive?}
The \vocab{contrapositive} of ``if $A$ then $B$'' is ``if not $B$ then not $A$''.
They are logically equivalent. Many proofs written as contradictions are really contrapositive
proofs: they assume not $B$, deduce not $A$ directly, and then say ``contradiction''.
Write those as contrapositive proofs. They are shorter and clearer.
A real contradiction proof uses the false assumption together with other facts and
reaches an impossibility in the middle of the argument, as in the three proofs above.

\section{Mathematical induction}

\subsection{The method}
\begin{theorem}[Induction principle]
  Let $P(n)$ be a statement depending on an integer $n$. Suppose
  \begin{enumerate}[(i)]
    \ii $P(n_0)$ is true (\vocab{base case}), and
    \ii for every $k \geq n_0$, if $P(k)$ is true then $P(k+1)$ is true (\vocab{inductive step}).
  \end{enumerate}
  Then $P(n)$ is true for all $n \geq n_0$.
\end{theorem}
The base case starts the chain and the step moves it forward one place at a time.
Both are required. In the step, $k$ is arbitrary: you may use only $P(k)$ and $k \geq n_0$.

\paragraph{When to use it.}
Use induction when the statement is indexed by an integer, or when the objects in the problem
can be \emph{built up} one piece at a time (add one line, one element, one vertex)
and the property is preserved at each piece. Most of the skill is in the inductive step.
The next three subsections give three ways to find the step.

\subsection{Find $P(k)$ inside $P(k+1)$}
The penultimate step for an inductive step is: \emph{write $P(k+1)$ and look for a copy of $P(k)$ in it.}
Whatever is left over is what you must handle.

\begin{theorem}[Fermat's little theorem]
  Let $p$ be a prime. Then $p \mid n^p - n$ for every positive integer $n$.
\end{theorem}

\paragraph{Investigation.}
$P(k)$ is ``$p \mid k^p - k$''. Write the $P(k+1)$ expression and expand with the binomial theorem:
\[
  (k+1)^p - (k+1) = \left( k^p + \sum_{j=1}^{p-1} \binom pj k^j + 1 \right) - k - 1
  = \underbrace{(k^p - k)}_{\text{copy of } P(k)} + \sum_{j=1}^{p-1} \binom pj k^j.
\]
The copy of $P(k)$ is divisible by $p$ by hypothesis. The leftover is a sum of terms
$\binom pj k^j$ with $1 \leq j \leq p-1$. So the step reduces to:
\emph{$p$ divides $\binom pj$ for $1 \le j \le p-1$.} Test $p = 5$: $\binom51, \dots, \binom54
= 5, 10, 10, 5$. All divisible by $5$. For a composite number, $\binom 42 = 6$ is not divisible by $4$,
so primality must be used.

\begin{proof}
First, for $1 \leq j \leq p-1$ we have $p! = \binom pj \, j! \, (p-j)!$.
The prime $p$ divides $p!$ but does not divide $j!\,(p-j)!$, because every factor in that product is
less than $p$. Hence $p \mid \binom pj$.

Now induct on $n$. Base case: $1^p - 1 = 0$ is divisible by $p$.
Inductive step: assume $p \mid k^p - k$. By the binomial theorem,
\[
  (k+1)^p - (k+1) = (k^p - k) + \sum_{j=1}^{p-1} \binom pj k^j,
\]
and every term on the right is divisible by $p$. So $p \mid (k+1)^p - (k+1)$.
\end{proof}

\subsection{Choose what to induct on}
Some problems contain no integer at all. Then you must decide what to induct on.
Get your hands dirty by building the object one piece at a time; the number of pieces is
usually the right variable.

\begin{example}[P\&B \S1.2; ACPS 2.3.7]
  Finitely many lines divide the plane into regions. Show that the regions can be colored with two
  colors so that any two regions sharing a boundary segment have different colors.
\end{example}

\paragraph{Investigation.}
Draw one line: two half-planes, one of each color. Draw a second line crossing it: four regions,
colored like a checkerboard. Now add a third line. It cuts some old regions in two,
and the two halves of each cut region have the same color, which is not allowed.

Wishful thinking: we want the two halves of every cut region to get different colors,
and we want old neighbors to stay different. Each cut region has one half on each side of the new line.
So: \emph{keep every color on one side of the new line and swap every color on the other side.}
Old neighbors on the same side are either both kept or both swapped, so they stay different.
Neighbors across the new line are halves of one old region, so after the swap they differ.
The variable to induct on is the number of lines.

\begin{proof}
Call a coloring \emph{good} if regions sharing a boundary segment have different colors.
We prove by induction on $n \geq 1$ that the regions formed by any $n$ lines have a good coloring.

Base case $n = 1$: color one half-plane black and the other white.

Inductive step: suppose any configuration of $k$ lines has a good coloring, and take $k+1$ lines.
Remove one line $\ell$, take a good coloring of the regions formed by the other $k$ lines,
then put $\ell$ back and reverse the colors of all regions on one side of $\ell$.
Let $R, R'$ be two new regions sharing a boundary segment $\sigma$.
If $\sigma$ is not on $\ell$, then $R$ and $R'$ lie on the same side of $\ell$ and lie in two
different old regions that shared a segment; those had different colors and both were reversed
or both were kept. If $\sigma$ is on $\ell$, then $R$ and $R'$ are the two parts of one old region,
and exactly one of them was reversed. In both cases $R$ and $R'$ have different colors.
\end{proof}

\subsection{Strengthen the statement}
Sometimes the inductive step fails, not because the statement is false,
but because $P(k)$ is too weak to give $P(k+1)$.
Then prove a \emph{stronger} statement. You need more to reach the new goal,
but you also get more to use in the step.

\begin{example}
  Prove that for every positive integer $n$,
  \[
    Q_n := \frac{1 \cdot 3 \cdot 5 \cdots (2n-1)}{2 \cdot 4 \cdot 6 \cdots (2n)} < \frac{1}{\sqrt{3n}}.
  \]
\end{example}

\paragraph{Investigation: the direct attempt.}
Note $Q_{n+1} = Q_n \cdot \frac{2n+1}{2n+2}$. Assume $Q_n < 1/\sqrt{3n}$. Then
\[
  Q_{n+1} < \frac{2n+1}{(2n+2)\sqrt{3n}},
\]
and we want this to be at most $1/\sqrt{3n+3}$. Squaring and clearing denominators, we need
\[
  (2n+1)^2(3n+3) \le (2n+2)^2 (3n),
  \quad\text{that is,}\quad
  12n^3 + 24n^2 + 15n + 3 \le 12n^3 + 24n^2 + 12n.
\]
This is false for every $n$. The step fails.

\paragraph{Investigation: wishful thinking.}
The bound $1/\sqrt{3n}$ seems too loose to carry forward. Try a family of tighter bounds,
$Q_n \le 1/\sqrt{3n + c}$ with a constant $c \geq 0$ to be chosen, and see what $c$ the step needs.
The same computation now requires
\[
  (2n+1)^2(3n+3+c) \le (2n+2)^2(3n+c).
\]
Expanding both sides, the $n^3$ and $n^2$ terms cancel, and the condition becomes
\[
  0 \le (4c - 3)\,n + 3(c - 1).
\]
With $c = 1$ the right side is $n \ge 0$, so the step works.
(With $c = 0$ it is $-3n - 3$, which is the failure above.)
The base case needs $Q_1 = \tfrac12 \le 1/\sqrt{3 + c}$, that is, $c \le 1$.
So $c = 1$ satisfies both requirements. Since $1/\sqrt{3n+1} < 1/\sqrt{3n}$,
the stronger statement implies the original one.

\begin{proof}
We prove $Q_n \le 1/\sqrt{3n+1}$ for all $n \geq 1$ by induction. This implies the claim,
since $1/\sqrt{3n+1} < 1/\sqrt{3n}$.

Base case: $Q_1 = \tfrac12 = 1/\sqrt{4}$.

Inductive step: assume $Q_k \le 1/\sqrt{3k+1}$. Then
$Q_{k+1} = Q_k \cdot \frac{2k+1}{2k+2} \le \frac{2k+1}{(2k+2)\sqrt{3k+1}}$,
so it suffices to show $\frac{2k+1}{(2k+2)\sqrt{3k+1}} \le \frac{1}{\sqrt{3k+4}}$,
which is equivalent to $(2k+1)^2(3k+4) \le (2k+2)^2(3k+1)$. Expanding,
\[
  12k^3 + 28k^2 + 19k + 4 \le 12k^3 + 28k^2 + 20k + 4,
\]
which holds since $k \ge 0$.
\end{proof}

\begin{remark}
  Method for finding the right strengthening: put a parameter in the statement
  (here $c$), run the inductive step with the parameter, and read off which values make the step
  and the base case both work. This is wishful thinking made precise.
  You will need this idea in this week's problem set.
\end{remark}

\subsection{A wrong induction}
\begin{example}[A false proof]
  ``Claim: in any set of $n \ge 1$ horses, all horses have the same color.
  Base case $n = 1$: true.
  Step: given $k+1$ horses $h_1, \dots, h_{k+1}$, the horses $h_1, \dots, h_k$ have one color
  and the horses $h_2, \dots, h_{k+1}$ have one color, by the hypothesis. These groups share $h_2$,
  so all $k+1$ horses have the same color.''
\end{example}
The error is in the step from $k = 1$ to $k = 2$. The two groups are $\{h_1\}$ and $\{h_2\}$,
and they share no horse. The step must work for \emph{every} $k \geq n_0$, including the smallest one.
\textbf{Always test your inductive step at $k = n_0$}, and look for any sentence in it that
silently assumes $k$ is large.

\end{document}
